%%%%%%%%%%%%%%%%%%%%%%%%%%%%%%%%%%%%%%%%%%%%%%%%%%%%%%%%%%%%%%%%%%%%%%%%%%%
% Función: capítulo de ejemplo de introducción de la tesis por compendio.
%
% Uso: copia de restauración; normalmente no debes modificarla.
% Tras restaurarla, modifica la copia activa para adaptarla a tu trabajo.
%%%%%%%%%%%%%%%%%%%%%%%%%%%%%%%%%%%%%%%%%%%%%%%%%%%%%%%%%%%%%%%%%%%%%%%%%%%

\chapter{Introducción}
\label{cha:compendium-introduction}

Como ejemplo concreto de presentación de un área de investigación, puedes contrastar la localización visual en espacios inteligentes~\cite{Sanabria2026Landmark}, la evaluación funcional en actividades cotidianas~\cite{nieva2024proposal} y la localización de barcos mediante sensado acústico distribuido~\cite{ramirez2026jstars}. Son ejemplos reales de líneas de investigación distintas, no publicaciones atribuidas al doctorando ficticio de esta plantilla.

\section{Contexto y motivación}
\subsection{Contexto de la investigación}

Utiliza este capítulo para presentar el contexto de la investigación, su motivación y el problema central abordado por la tesis. Explica por qué el problema es relevante, delimita el alcance del trabajo y proporciona la información necesaria para comprender el resumen extendido. Este ejemplo también muestra el uso de acrónimos: \gls{AI} y \gls{HMI}.

\subsection{Planteamiento del problema}

Define el problema concreto abordado por la tesis, su alcance y las hipótesis que delimitan la investigación propuesta.

\section{Organización de la tesis}

Describe brevemente el contenido del resumen extendido y explica cómo las publicaciones incluidas en la parte~\ref{part:compendium-publications} sustentan el hilo argumental de la tesis.

%%% Local Variables:
%%% TeX-master: "../../book"
%%% End:
